% 第九卷 3.1.5 节：拼接序列上的普通因果掩码与文档内掩码。
% 序列长度 8：位置 1--3 属于文档甲，位置 4--8 属于文档乙。
% 行为查询位置 i，列为键位置 j；蓝格为同一文档内可见，黄格为普通因果掩码下的跨文档可见，白格不可见。下方一行标出两种做法的位置编号。
\begin{tikzpicture}[
  x=1cm,
  y=1cm,
  font=\figurefont\footnotesize,
  seen/.style={draw=white,line width=0.5pt,fill=FigureBlue!70},
  cross/.style={draw=white,line width=0.5pt,fill=FigureYellow!85},
  hidden/.style={draw=FigureGrid,line width=0.4pt,fill=white},
  idx/.style={font=\figurefont\scriptsize,text=FigureMuted},
  ptitle/.style={font=\figurefont\bfseries\footnotesize,text=FigureInk,anchor=south},
  note/.style={font=\figurefont\scriptsize,text=FigureMuted,anchor=north,align=center}
]
  \def\s{0.34}
  % 面板编号 \p：1 普通因果掩码，2 文档内掩码。
  % 普通因果掩码中跨文档的可见格用黄色标出，便于与文档内可见格区分。
  \foreach \p/\xo/\ttl in {%
      1/0.0/{(a) 普通因果掩码},%
      2/4.3/{(b) 文档内掩码}} {
    \foreach \i in {1,...,8} {
      \foreach \j in {1,...,8} {
        \pgfmathtruncatemacro{\kind}{%
          \j>\i ? 0 :
          ((\i<=3 || \j>=4) ? 1 : (\p==1 ? 2 : 0))}
        \ifcase\kind
          \path[hidden] ({\xo+(\j-1)*\s},{-(\i-1)*\s}) rectangle ++(\s,-\s);
        \or
          \path[seen] ({\xo+(\j-1)*\s},{-(\i-1)*\s}) rectangle ++(\s,-\s);
        \or
          \path[cross] ({\xo+(\j-1)*\s},{-(\i-1)*\s}) rectangle ++(\s,-\s);
        \fi
      }
      \node[idx,anchor=east] at ({\xo-0.04},{-(\i-0.5)*\s}) {\i};
      \node[idx,anchor=south] at ({\xo+(\i-0.5)*\s},0.02) {\i};
    }
    \draw[FigureStroke,line width=0.6pt] (\xo,0) rectangle ++({8*\s},{-8*\s});
    % 文档边界。
    \draw[FigureStroke,line width=0.9pt,dashed] ({\xo+3*\s},0) -- ({\xo+3*\s},{-8*\s});
    \draw[FigureStroke,line width=0.9pt,dashed] (\xo,{-3*\s}) -- ({\xo+8*\s},{-3*\s});
    \node[ptitle] at ({\xo+4*\s},0.32) {\ttl};
  }
  % 位置编号。
  \node[note] at ({4*\s},{-8*\s-0.12}) {位置编号：1\enspace2\enspace3\enspace4\enspace5\enspace6\enspace7\enspace8};
  \node[note] at ({4.3+4*\s},{-8*\s-0.12}) {位置编号：1\enspace2\enspace3\quad1\enspace2\enspace3\enspace4\enspace5};
  % 坐标说明与图例。
  \node[idx,anchor=east] at (-0.3,{-4*\s}) {查询$i$};
  \node[idx,anchor=south east] at (-0.06,0.02) {键$j$};
  \path[seen] (0.2,{-8*\s-0.75}) rectangle ++(\s,-\s);
  \node[anchor=west,font=\figurefont\scriptsize,text=FigureInk] at ({0.2+\s+0.08},{-8*\s-0.75-0.5*\s}) {同一文档内可见};
  \path[cross] (2.9,{-8*\s-0.75}) rectangle ++(\s,-\s);
  \node[anchor=west,font=\figurefont\scriptsize,text=FigureInk] at ({2.9+\s+0.08},{-8*\s-0.75-0.5*\s}) {跨文档可见};
  \path[hidden] (5.3,{-8*\s-0.75}) rectangle ++(\s,-\s);
  \node[anchor=west,font=\figurefont\scriptsize,text=FigureInk] at ({5.3+\s+0.08},{-8*\s-0.75-0.5*\s}) {不可见};
\end{tikzpicture}
